\documentclass[10pt]{article}

\usepackage[left=0.7in,right=0.7in,top=0.35in,bottom=0.7in]{geometry}
\usepackage{xcolor}
\usepackage{graphicx}
\usepackage{hyperref}

\hypersetup{colorlinks=true,linkcolor=blue,urlcolor=blue}
\urlstyle{rm}
\usepackage{url}


\title{CAI 6108: Machine Learning Engineering\\
	\large Project Report:  {\textcolor{purple}{Group 41}}} %% TODO: replace with the title of your project
	
	
	
%% TODO: your name and email go here (all members of the group)
%% Comment out as needed and designate a point of contact
%% Add / remove names as necessary
\author{
        Addison Ferrell \\
        aferrell2@ufl.edu\\
        \and
        Sai Pranav Reddy Guduru \\
        saipranavrguduru@ufl.edu\\
        \and
        Ametoje Virdi \\
        ametojevirdi@ufl.edu\\
        \and
        William Le \\
        lew@ufl.edu\\
}


% set the date to today
\date{April 24, 2026}


\begin{document} % start document tag

\maketitle



%%% Remember: writing counts! (try to be clear and concise.)
%%% Make sure to cite your sources and references (use refs.bib and \cite{} or \footnote{} for URLs).
%%%



%% TODO: write an introduction to make the report self-contained
%% Must address:
%% - What is the project about and what is notable about the approach and results?
%%
\section{Introduction}

This course project is about building a complete machine learning pipeline for
multilabel image classification. The task is to label a $128\times128$ PNG
image with any subset of 12 object classes: pen, paper, book, clock, phone,
laptop, chair, desk, bottle, keychain, backpack, and calculator. Because an
image can contain multiple objects at once, this is a multilabel classification
problem rather than a standard single-label classification task.

Our goal is to make accurate predictions on unseen data while following the
project constraints from class: the pipeline must run entirely offline, use
course-relevant Python libraries, and apply reasonable machine learning best
practices. We therefore built a PyTorch-based training and evaluation pipeline
that supports multiple pretrained vision backbones and compares different
fine-tuning strategies across them.

The main idea behind our experiments is to test whether full fine-tuning of a
pretrained model works better than training only the final classification head,
and whether that trend is consistent across several architectures. The models
we compare are ResNet-18, DenseNet-121, EfficientNet-B0, and
ViT-B/32~\cite{he2016deep,huang2017densely,tan2019efficientnet,dosovitskiy2021image}.
The notable part of the project is not just the final accuracy, but also the
fact that the same offline pipeline, shared data split, and common metric set
are used across all experiments so the results can be compared fairly.




%% TODO: write about your approach / ML pipeline
%% Must contain:
%% - How you are trying to solve this problem
%% - How did you process the data?
%% - What is the task and approach (ML techniques)?
%%
\section{Approach}

Our pipeline is implemented entirely in Python using PyTorch and torchvision,
which keeps the project offline and within the toolset used in the course. The
dataset is organized into folders whose names encode one or more labels using
underscores. For example, a folder such as
\texttt{phone\_laptop\_keychain} indicates that every image in that folder
contains those labels. The loader converts each image into a 12-dimensional
multi-hot target vector using a fixed class order shared by training,
checkpoint saving, and evaluation.

For preprocessing, every image is converted to RGB, resized, transformed to a
tensor, and normalized using ImageNet mean and standard deviation. During
training we apply light data augmentation with random horizontal flips and
small color jitter to improve robustness. Most architectures use
$128\times128$ inputs, while ViT-B/32 uses $224\times224$ inputs because that
matches its expected resolution.

The model pipeline supports four backbones: \texttt{resnet18},
\texttt{densenet121}, \texttt{efficientnet\_b0}, and \texttt{vitb32}. In each
case, we replace the original classifier with a 12-output layer so the network
produces one logit per class. Since the task is multilabel, we train with
\texttt{BCEWithLogitsLoss}. At inference time, we apply a sigmoid to convert
logits to probabilities and then threshold the outputs at 0.5 to decide which
labels are predicted.

We compare two training pipelines for each architecture. In full fine-tuning,
all pretrained parameters are updated on our dataset. In head-only
classification, the backbone is frozen and only the final classifier layer is
trained. We use pretrained torchvision weights when available, optimize with
AdamW, and save the best checkpoint according to validation micro-F1. This lets
us study both architecture choice and transfer-learning strategy while keeping
the rest of the pipeline fixed.





%% TODO: write about your evaluation methodology
%% Must contain:
%% - What are the metrics?
%% - What are the baselines?
%% - How did you split the data?
%%
\section{Evaluation Methodology}

We follow the class recommendation to use a train/validation/test split and to
evaluate on unseen data. Our experiments use a fixed 70\%/15\%/15\% split with
random seed 42, stored in \texttt{splits/split\_seed42.json}. Using the same
split for all runs makes the comparisons fair across architectures, learning
rates, and fine-tuning strategies. The train split is used for optimization,
the validation split is used for model selection, and the test split is used
only for final reporting.

The pipeline records several evaluation metrics, including exact-match
accuracy, hamming accuracy, mean intersection over union (IoU), micro
precision, micro recall, micro-F1, and macro-F1. We use validation micro-F1 as
the main model-selection metric because it balances precision and recall across
all predicted labels and is appropriate for multilabel classification.
Exact-match accuracy is stricter because every label for an image must be
correct, while hamming accuracy gives a more label-wise view of performance.
Macro-F1 is useful because it gives equal importance to each class.

Our main baselines are the simpler versions of the same pipeline. For each
architecture, head-only training acts as a baseline for full fine-tuning. We
also compare different learning rates within the same architecture. This means
the Results section should answer the questions emphasized in the project
slides: what results we obtained, how those results compare to the baselines,
and whether the pipeline appears to follow good machine learning practice.

One limitation of our setup is that the split is created at the image level
rather than the folder level. Because of that, related images from the same
object combination may appear in different partitions, which can make the task
slightly easier than a stricter split that holds out entire folders.




%% TODO: write about your results/findings
%% Must contain:
%% - results comparing your approach to baseline according to metrics
%%
%% You can present this information in whatever way you want but consider tables and (or) figures.
%%
\section{Results}

Table~\ref{tab:results} reports all test-set metrics across all 16 configurations.
The best value in each column is \textbf{bold}.

\begin{table}[h]
\centering
\caption{Test-set results. Full FT = full fine-tuning; Head = head-only (backbone frozen).
Prec.\ = micro precision, Rec.\ = micro recall.}
\label{tab:results}
\resizebox{\textwidth}{!}{%
\begin{tabular}{llrrrrrrr}
\hline
\textbf{Model} & \textbf{Strategy} & \textbf{LR} & \textbf{Exact Match} & \textbf{Hamming} & \textbf{IoU} & \textbf{Prec.} & \textbf{Rec.} & \textbf{F1\textsubscript{micro}} \\
\hline
DenseNet-121    & Full FT & 1e-4 & \textbf{0.669} & \textbf{0.958} & \textbf{0.811} & \textbf{0.881} & \textbf{0.817} & \textbf{0.848} \\
DenseNet-121    & Full FT & 3e-4 & 0.614 & 0.948 & 0.766 & 0.847 & 0.775 & 0.810 \\
DenseNet-121    & Head    & 1e-3 & 0.362 & 0.913 & 0.584 & 0.754 & 0.586 & 0.660 \\
DenseNet-121    & Head    & 3e-4 & 0.372 & 0.914 & 0.587 & 0.775 & 0.566 & 0.655 \\
\hline
EfficientNet-B0 & Full FT & 1e-4 & 0.606 & 0.949 & 0.775 & 0.850 & 0.783 & 0.815 \\
EfficientNet-B0 & Full FT & 3e-4 & 0.597 & 0.946 & 0.763 & 0.841 & 0.770 & 0.804 \\
EfficientNet-B0 & Head    & 1e-3 & 0.312 & 0.903 & 0.513 & 0.720 & 0.535 & 0.614 \\
EfficientNet-B0 & Head    & 3e-4 & 0.261 & 0.892 & 0.461 & 0.674 & 0.484 & 0.563 \\
\hline
ResNet-18       & Full FT & 1e-4 & 0.543 & 0.942 & 0.723 & 0.849 & 0.724 & 0.782 \\
ResNet-18       & Full FT & 3e-4 & 0.531 & 0.939 & 0.709 & 0.830 & 0.722 & 0.772 \\
ResNet-18       & Head    & 1e-3 & 0.324 & 0.907 & 0.535 & 0.749 & 0.529 & 0.620 \\
ResNet-18       & Head    & 3e-4 & 0.311 & 0.904 & 0.499 & 0.768 & 0.473 & 0.585 \\
\hline
ViT-B/32        & Full FT & 1e-4 & 0.610 & 0.946 & 0.769 & 0.821 & 0.796 & 0.808 \\
ViT-B/32        & Full FT & 3e-4 & 0.195 & 0.842 & 0.320 & 0.437 & 0.354 & 0.392 \\
ViT-B/32        & Head    & 1e-3 & 0.503 & 0.934 & 0.706 & 0.814 & 0.699 & 0.752 \\
ViT-B/32        & Head    & 3e-4 & 0.485 & 0.934 & 0.696 & 0.829 & 0.681 & 0.747 \\
\hline
\end{tabular}}
\end{table}

\textbf{Best model.}
DenseNet-121 with full fine-tuning at $\text{lr}=10^{-4}$ achieves the best score
on every metric as it yielded a 66.9\% exact match, 95.8\% hamming accuracy, 81.1\% mean IoU,
88.1\% micro precision, 81.7\% micro recall, and 84.8\% micro-F1.

\textbf{Full fine-tuning vs.\ head-only.}
Full fine-tuning outperformed head-only training for all three CNN architectures, 
as evidenced by the 18--25 point gap in micro-F1. Recall is primarily responsible for
this gap since head-only models have reasonable precision but cannot adapt their features
to the dataset, resulting in many missed true labels. Full fine-tuning allows the parameters
to adjust to the dataset, which significantly improves recall and closes the F1 gap.

\textbf{ViT behaves differently.}
ViT-B/32 head-only reaches 75.2\% micro-F1, which is well above the CNN head-only
results. However, ViT's full fine-tuning micro-f1 at $\text{lr}=3\times10^{-4}$ falls to
39.2\%, with both precision (0.437) and recall (0.354) significantly dropping.
The model clearly experiences training instability at the higher learning rate.
Lowering the learning rate to $10^{-4}$ increases the micro-f1 to 80.8\%, but ViT still does not
beat DenseNet-121 even while having 87M parameters as compared with DenseNet's 7M.

\textbf{Learning rate sensitivity.}
For CNNs, the two learning rates typically differ by no more than 10 points across all metrics,
with $10^{-4}$ consistently better. ViT is an outlier: the higher rate causes a
40-point collapse in micro-F1 under full fine-tuning but has negligible effect
in head-only mode, where only 9k parameters are updated.

\textbf{Precision vs.\ recall tradeoff.}
Across all configurations, precision is consistently higher than recall. Models
tend to be conservative, predicting fewer labels than are present, rather than
over-predicting. Full fine-tuning improves recall far more than precision,
which is why it closes the F1 gap so significantly over head-only training.

%% TODO: write about what you conclude. This is not meant to be a summary section but more of a takeaways/future work section.
%% Must contain:
%% - Any concrete takeaways or conclusions from your experiments/results/project
%%
%% You can also discuss limitations here if you want.
%%
\section{Conclusions}

This section should summarize the main takeaways after the final runs are
aggregated. Good points to cover:

\begin{itemize}
    \item the best overall model and training setup
    \item whether full fine-tuning was worth the extra compute
    \item how architecture choice affected multilabel performance
    \item limitations of the current dataset split and dataset size
    \item one or two concrete next steps, such as stronger augmentation,
    threshold tuning, or a folder-level split
    \item whether the final pipeline appears reasonable for unseen offline evaluation
\end{itemize}

%%%%

\bibliography{refs}
\bibliographystyle{plain}


\end{document} % end tag of the document
